\begin{figure}[tbp]
\centering
\setlength{\unitlength}{1pt}
\begin{picture}(430,145)
  \put(100,125){\makebox(0,0){\small normal real slice ($v=0$)}}
  \put(15,75){\vector(1,0){180}}
  {\color{blue!65!black}\linethickness{2pt}\put(35,75){\line(1,0){130}}}\ignorespaces
  \put(100,71){\line(0,1){8}}
  \put(116,68){\line(0,1){14}}
  \put(100,75){\circle*{3}}
  \put(116,75){\circle*{3}}
  \put(91,58){\makebox(0,0){\small $x_j=-\delta_j$}}
  \put(127,91){\makebox(0,0){\small $p: u=0$}}
  \put(35,101){\line(1,0){130}}
  \put(35,98){\line(0,1){6}}
  \put(165,98){\line(0,1){6}}
  \put(100,111){\makebox(0,0){\small shell width $2\eta\lambda_j$}}
  \put(193,63){\makebox(0,0){\small $u$}}
  \put(108,37){\makebox(0,0){\small $\delta_j/\lambda_j\longrightarrow0$}}
  \put(320,125){\makebox(0,0){\small complex tangent slice ($u=x_j$)}}
  \put(320,69){\circle{90}}
  {\color{blue!65!black}\put(320,69){\circle{34}}\put(320,69){\circle{43}}}\ignorespaces
  \put(320,69){\circle*{2}}
  \put(320,69){\vector(1,0){45}}
  \put(376,58){\makebox(0,0)[l]{\small $r(x_j)$}}
  \put(320,69){\vector(-4,3){16}}
  \put(300,90){\makebox(0,0)[r]{\small $\lambda_j$}}
  \put(320,8){\makebox(0,0){\small $\lambda_j/\sqrt{\delta_j}\longrightarrow0$}}
\end{picture}
\caption{The two scale inequalities used to cross the boundary. On the
left, the filled Hartogs domain extends in the $u$ direction farther than
the distance $\delta_j$ from the slice to $p$. On the right, its shell
has radius comparable to $\lambda_j$ and fits well inside the slice disk
of radius $r(x_j)=k\sqrt{\delta_j}$. The annulus has a small image in one
target product chart. These are separate source slices in the affine
coordinates $(u,v)$, drawn schematically and not at a common scale.}
\label{fig:boundary-scales}
\end{figure}
